% !TEX program = xelatex
% =====================================================================
%  PhotoLib 管理员使用指南
%  编译：xelatex 管理员使用指南.tex（连续两次以生成目录/交叉引用）
% =====================================================================
\documentclass[cn,11pt,green]{elegantbook}

% =====================================================================
%  elegant-style.tex —— PhotoLib 使用指南三份文档共用的样式与宏定义
%  由 管理员/部长/校区负责人 三份指南通过 % =====================================================================
%  elegant-style.tex —— PhotoLib 使用指南三份文档共用的样式与宏定义
%  由 管理员/部长/校区负责人 三份指南通过 % =====================================================================
%  elegant-style.tex —— PhotoLib 使用指南三份文档共用的样式与宏定义
%  由 管理员/部长/校区负责人 三份指南通过 \input{elegant-style} 引入。
%  依赖：ElegantBook 文档类（须以 XeLaTeX 编译）。
% =====================================================================

% ---- 颜色（与 ElegantBook green 主题协调）----
\definecolor{plMain}{RGB}{0,120,80}      % 主色：墨绿
\definecolor{plInk}{RGB}{38,50,56}       % 正文强调
\definecolor{admTipBg}{RGB}{236,247,240}
\definecolor{admTipFr}{RGB}{34,139,88}
\definecolor{admNoteBg}{RGB}{232,240,254}
\definecolor{admNoteFr}{RGB}{25,103,210}
\definecolor{admWarnBg}{RGB}{253,237,237}
\definecolor{admWarnFr}{RGB}{198,40,40}
\definecolor{admStepBg}{RGB}{245,243,255}
\definecolor{admStepFr}{RGB}{103,58,183}

% ---- tcolorbox（ElegantBook 已载入，这里补声明所需库）----
\usepackage{tcolorbox}
\tcbuselibrary{skins,breakable}

% 统一的四种“提示框”环境，均支持可选自定义标题：
%   \begin{tipbox}[自定义标题] ... \end{tipbox}
\newtcolorbox{tipbox}[1][小贴士]{%
  enhanced, breakable,
  colback=admTipBg, colframe=admTipFr, coltitle=white,
  fonttitle=\bfseries\sffamily, title={#1},
  boxrule=0.6pt, arc=3pt, left=8pt, right=8pt, top=6pt, bottom=6pt}

\newtcolorbox{notebox}[1][说明]{%
  enhanced, breakable,
  colback=admNoteBg, colframe=admNoteFr, coltitle=white,
  fonttitle=\bfseries\sffamily, title={#1},
  boxrule=0.6pt, arc=3pt, left=8pt, right=8pt, top=6pt, bottom=6pt}

\newtcolorbox{warnbox}[1][注意]{%
  enhanced, breakable,
  colback=admWarnBg, colframe=admWarnFr, coltitle=white,
  fonttitle=\bfseries\sffamily, title={#1},
  boxrule=0.6pt, arc=3pt, left=8pt, right=8pt, top=6pt, bottom=6pt}

\newtcolorbox{stepbox}[1][操作步骤]{%
  enhanced, breakable,
  colback=admStepBg, colframe=admStepFr, coltitle=white,
  fonttitle=\bfseries\sffamily, title={#1},
  boxrule=0.6pt, arc=3pt, left=8pt, right=8pt, top=6pt, bottom=6pt}

% ---- 表格 ----
\usepackage{booktabs}
\usepackage{array}
\usepackage{longtable}
\usepackage{makecell}
\renewcommand{\arraystretch}{1.28}

% 左对齐、可换行的定宽列
\newcolumntype{L}[1]{>{\raggedright\arraybackslash}p{#1}}

% ---- 行内“键位/菜单/字段”标记 ----
\usepackage{relsize}
% 菜单路径：系统管理 \menu{数据备份}
\newcommand{\menu}[1]{{\sffamily\bfseries\color{plMain}#1}}
% 界面按钮/字段名
\newcommand{\ui}[1]{\textsf{「#1」}}
% 权限码等等宽标识；在每个下划线后允许换行，避免长标识顶出表格边界
\newcommand{\code}[1]{{\ttfamily\smaller[0.5]%
  \renewcommand{\_}{\textunderscore\penalty0\hspace{0pt}}#1}}
% 角色徽标
\newcommand{\role}[1]{{\sffamily\bfseries\color{admNoteFr}#1}}

% ---- 超链接 ----
\hypersetup{
  colorlinks=true,
  linkcolor=plMain,
  urlcolor=admNoteFr,
  citecolor=plMain}
 引入。
%  依赖：ElegantBook 文档类（须以 XeLaTeX 编译）。
% =====================================================================

% ---- 颜色（与 ElegantBook green 主题协调）----
\definecolor{plMain}{RGB}{0,120,80}      % 主色：墨绿
\definecolor{plInk}{RGB}{38,50,56}       % 正文强调
\definecolor{admTipBg}{RGB}{236,247,240}
\definecolor{admTipFr}{RGB}{34,139,88}
\definecolor{admNoteBg}{RGB}{232,240,254}
\definecolor{admNoteFr}{RGB}{25,103,210}
\definecolor{admWarnBg}{RGB}{253,237,237}
\definecolor{admWarnFr}{RGB}{198,40,40}
\definecolor{admStepBg}{RGB}{245,243,255}
\definecolor{admStepFr}{RGB}{103,58,183}

% ---- tcolorbox（ElegantBook 已载入，这里补声明所需库）----
\usepackage{tcolorbox}
\tcbuselibrary{skins,breakable}

% 统一的四种“提示框”环境，均支持可选自定义标题：
%   \begin{tipbox}[自定义标题] ... \end{tipbox}
\newtcolorbox{tipbox}[1][小贴士]{%
  enhanced, breakable,
  colback=admTipBg, colframe=admTipFr, coltitle=white,
  fonttitle=\bfseries\sffamily, title={#1},
  boxrule=0.6pt, arc=3pt, left=8pt, right=8pt, top=6pt, bottom=6pt}

\newtcolorbox{notebox}[1][说明]{%
  enhanced, breakable,
  colback=admNoteBg, colframe=admNoteFr, coltitle=white,
  fonttitle=\bfseries\sffamily, title={#1},
  boxrule=0.6pt, arc=3pt, left=8pt, right=8pt, top=6pt, bottom=6pt}

\newtcolorbox{warnbox}[1][注意]{%
  enhanced, breakable,
  colback=admWarnBg, colframe=admWarnFr, coltitle=white,
  fonttitle=\bfseries\sffamily, title={#1},
  boxrule=0.6pt, arc=3pt, left=8pt, right=8pt, top=6pt, bottom=6pt}

\newtcolorbox{stepbox}[1][操作步骤]{%
  enhanced, breakable,
  colback=admStepBg, colframe=admStepFr, coltitle=white,
  fonttitle=\bfseries\sffamily, title={#1},
  boxrule=0.6pt, arc=3pt, left=8pt, right=8pt, top=6pt, bottom=6pt}

% ---- 表格 ----
\usepackage{booktabs}
\usepackage{array}
\usepackage{longtable}
\usepackage{makecell}
\renewcommand{\arraystretch}{1.28}

% 左对齐、可换行的定宽列
\newcolumntype{L}[1]{>{\raggedright\arraybackslash}p{#1}}

% ---- 行内“键位/菜单/字段”标记 ----
\usepackage{relsize}
% 菜单路径：系统管理 \menu{数据备份}
\newcommand{\menu}[1]{{\sffamily\bfseries\color{plMain}#1}}
% 界面按钮/字段名
\newcommand{\ui}[1]{\textsf{「#1」}}
% 权限码等等宽标识；在每个下划线后允许换行，避免长标识顶出表格边界
\newcommand{\code}[1]{{\ttfamily\smaller[0.5]%
  \renewcommand{\_}{\textunderscore\penalty0\hspace{0pt}}#1}}
% 角色徽标
\newcommand{\role}[1]{{\sffamily\bfseries\color{admNoteFr}#1}}

% ---- 超链接 ----
\hypersetup{
  colorlinks=true,
  linkcolor=plMain,
  urlcolor=admNoteFr,
  citecolor=plMain}
 引入。
%  依赖：ElegantBook 文档类（须以 XeLaTeX 编译）。
% =====================================================================

% ---- 颜色（与 ElegantBook green 主题协调）----
\definecolor{plMain}{RGB}{0,120,80}      % 主色：墨绿
\definecolor{plInk}{RGB}{38,50,56}       % 正文强调
\definecolor{admTipBg}{RGB}{236,247,240}
\definecolor{admTipFr}{RGB}{34,139,88}
\definecolor{admNoteBg}{RGB}{232,240,254}
\definecolor{admNoteFr}{RGB}{25,103,210}
\definecolor{admWarnBg}{RGB}{253,237,237}
\definecolor{admWarnFr}{RGB}{198,40,40}
\definecolor{admStepBg}{RGB}{245,243,255}
\definecolor{admStepFr}{RGB}{103,58,183}

% ---- tcolorbox（ElegantBook 已载入，这里补声明所需库）----
\usepackage{tcolorbox}
\tcbuselibrary{skins,breakable}

% 统一的四种“提示框”环境，均支持可选自定义标题：
%   \begin{tipbox}[自定义标题] ... \end{tipbox}
\newtcolorbox{tipbox}[1][小贴士]{%
  enhanced, breakable,
  colback=admTipBg, colframe=admTipFr, coltitle=white,
  fonttitle=\bfseries\sffamily, title={#1},
  boxrule=0.6pt, arc=3pt, left=8pt, right=8pt, top=6pt, bottom=6pt}

\newtcolorbox{notebox}[1][说明]{%
  enhanced, breakable,
  colback=admNoteBg, colframe=admNoteFr, coltitle=white,
  fonttitle=\bfseries\sffamily, title={#1},
  boxrule=0.6pt, arc=3pt, left=8pt, right=8pt, top=6pt, bottom=6pt}

\newtcolorbox{warnbox}[1][注意]{%
  enhanced, breakable,
  colback=admWarnBg, colframe=admWarnFr, coltitle=white,
  fonttitle=\bfseries\sffamily, title={#1},
  boxrule=0.6pt, arc=3pt, left=8pt, right=8pt, top=6pt, bottom=6pt}

\newtcolorbox{stepbox}[1][操作步骤]{%
  enhanced, breakable,
  colback=admStepBg, colframe=admStepFr, coltitle=white,
  fonttitle=\bfseries\sffamily, title={#1},
  boxrule=0.6pt, arc=3pt, left=8pt, right=8pt, top=6pt, bottom=6pt}

% ---- 表格 ----
\usepackage{booktabs}
\usepackage{array}
\usepackage{longtable}
\usepackage{makecell}
\renewcommand{\arraystretch}{1.28}

% 左对齐、可换行的定宽列
\newcolumntype{L}[1]{>{\raggedright\arraybackslash}p{#1}}

% ---- 行内“键位/菜单/字段”标记 ----
\usepackage{relsize}
% 菜单路径：系统管理 \menu{数据备份}
\newcommand{\menu}[1]{{\sffamily\bfseries\color{plMain}#1}}
% 界面按钮/字段名
\newcommand{\ui}[1]{\textsf{「#1」}}
% 权限码等等宽标识；在每个下划线后允许换行，避免长标识顶出表格边界
\newcommand{\code}[1]{{\ttfamily\smaller[0.5]%
  \renewcommand{\_}{\textunderscore\penalty0\hspace{0pt}}#1}}
% 角色徽标
\newcommand{\role}[1]{{\sffamily\bfseries\color{admNoteFr}#1}}

% ---- 超链接 ----
\hypersetup{
  colorlinks=true,
  linkcolor=plMain,
  urlcolor=admNoteFr,
  citecolor=plMain}


\title{PhotoLib 管理员使用指南}
\subtitle{系统管理 · 技术架构 · 部署运维}
\author{PhotoLib 项目组}
\institute{校园摄影部图片工作站}
\version{1.0}
\date{2026 年 9 月}
\extrainfo{本指南面向系统管理员（ADMIN），介绍系统的技术架构、部署方式、\\后台管理功能与日常运维要点。}
\cover{cover.png}

\begin{document}

\maketitle
\frontmatter
\tableofcontents
\mainmatter

% =====================================================================
\chapter{认识这套系统}
% =====================================================================

\begin{introduction}
  \item PhotoLib 是什么
  \item 它由哪些部分组成
  \item 管理员在其中扮演什么角色
\end{introduction}

\section{一句话介绍}

PhotoLib 是一套面向校园摄影部门的\textbf{图片工作站}。它把摄影部日常的一整条工作线串了起来：
立项选题 → 发布图片需求 → 校区负责人接单 → 上传与管理图片 → 标记采用 → 统计被引张数和工时 → 导出数据。
此外还带有站内消息、成员招募、好图精选、文档中心等协作功能。

系统只有一个网站入口，所有人（管理员、部长、校区负责人）都用同一个地址登录，看到的菜单和能做的操作由\textbf{账号权限}决定。

\section{系统由哪几块组成}

\begin{table}[htbp]
  \centering
  \caption{系统的技术组成}
  \begin{tabular}{L{2.6cm} L{9.5cm}}
    \toprule
    \textbf{部分} & \textbf{说明} \\
    \midrule
    前端页面 & 用户在浏览器里看到的界面。用 React 19 + Ant Design 6 编写，单页应用。 \\
    后端服务 & 处理业务逻辑、权限校验、数据读写。用 Java 21 + Spring Boot 4 编写。 \\
    数据库 & MySQL 8，保存账号、项目、图片记录、工时等结构化数据。 \\
    文件存储 & 保存图片本身。生产环境用阿里云 OSS（私有桶），本地开发用磁盘。 \\
    邮件服务 & 阿里云 DirectMail，用于发送通知邮件。 \\
    图片处理 & 一套原生组件（Zig + libjpeg-turbo/stb/libvips），负责压缩、生成预览图。 \\
    \bottomrule
  \end{tabular}
\end{table}

\begin{notebox}[打包方式]
  前端和后端\textbf{最终打包成同一个 JAR 文件}运行。也就是说，部署时你只需要一个 Java 程序，
  不需要单独部署前端网站、也不需要在服务器上装 Node。构建时 Maven 会自动把前端页面塞进后端 JAR 的静态资源里。
\end{notebox}

\section{管理员的定位}

管理员（角色代码 \code{ADMIN}）拥有系统\textbf{最高权限}，可以管理一切资源：账号、校区、权限组、品牌、
审计日志、存储对账、数据库备份等。本指南的重点，就是帮助你理解系统结构，并用好这些管理能力。

\begin{warnbox}[一条贯穿全书的安全原则]
  \textbf{后端才是权限的最终边界。} 前端「藏起某个按钮」只是界面效果，不等于真的没有权限。
  任何一次操作能不能做，都由后端根据账号权限重新判断。因此配置权限时，请以\textbf{权限组的授权}为准，
  不要以为「界面上看不到就做不了」。
\end{warnbox}

% =====================================================================
\chapter{技术架构详解}
% =====================================================================

\begin{introduction}
  \item 前后端如何分工
  \item 图片是怎么被处理和存储的
  \item 系统靠哪些机制保证安全与稳定
\end{introduction}

\section{整体结构}

所有后端接口都以 \code{/api/v1} 开头。前端使用\textbf{哈希路由}（地址里带 \code{\#}，例如 \code{/\#/projects}），
这样即使前端被打包进后端 JAR、从服务器根路径访问，刷新页面也不会 404。

后端按业务领域分包，每个包大致遵循三层结构：

\begin{itemize}
  \item \textbf{Controller}：接收 HTTP 请求、校验参数、做权限判断。
  \item \textbf{Service}：真正的业务规则都写在这里。
  \item \textbf{Mapper}：用 MyBatis-Plus 访问数据库。
\end{itemize}

主要业务包包括：项目（project）、需求（request）、图片（photo）、采用（adoption）、工时（worklog）、
统计（statistics）、通知（notification）、审计（audit）、存储（storage）、招募（recruitment）、
文档（doc）、迁移（migration）、认证（auth）、后台（admin）等。

\section{图片是怎么处理的}

图片从上传到能在图库里看到，中间经过一条受控的处理流水线：

\begin{enumerate}
  \item 浏览器先算出图片的 SHA-256 指纹，向后端申请一个\textbf{短时签名的直传地址}。
  \item 浏览器把原图\textbf{直接传到对象存储}（不经过后端服务器中转）。
  \item 后端重新读取原图，校验文件类型、指纹是否一致，然后压缩出「成品图」和「预览图」。
  \item 成品图用于下载，预览图（统一为 WebP 格式）用于图库缩略图展示。
\end{enumerate}

\begin{notebox}[为什么预览图用 WebP]
  同样的清晰度下，WebP 预览比 JPEG 小约三成、比无损 PNG 小十几倍，还自带透明通道支持。
  预览图质量由 \code{PREVIEW\_COMPRESSION\_RATIO} 控制（取值 0 到 1，默认 0.6）。
\end{notebox}

\begin{table}[htbp]
  \centering
  \caption{图片相关的关键限制}
  \begin{tabular}{L{5.5cm} L{6.6cm}}
    \toprule
    \textbf{项目} & \textbf{限制} \\
    \midrule
    支持的格式 & 仅 JPG / JPEG / PNG \\
    单张图片大小 & 最大 100 MiB \\
    成品图目标体积 & 10 MiB \\
    单次批量上传 & 最多 100 张 \\
    ZIP 上传 & 压缩包最大约 1.5 GB，解压后总量最大 10 GiB \\
    原图保留期 & 默认 30 天（生成成品图后才会按期清理原图） \\
    签名下载地址有效期 & 默认 15 分钟 \\
    \bottomrule
  \end{tabular}
\end{table}

\section{系统靠什么保证安全和稳定}

\begin{itemize}
  \item \textbf{登录安全}：访问令牌 + HttpOnly 刷新令牌；首次登录强制改密；登录失败限速与锁定；登录尝试全量审计。
  \item \textbf{私有存储}：图片放在私有对象存储桶里，浏览器只能拿到\textbf{短时签名地址}，不会暴露永久公共链接。
  \item \textbf{并发保护}：关键写操作使用乐观锁（版本号）和幂等键，避免重复提交、并发覆盖。
  \item \textbf{软删除}：图片、需求等\textbf{不做物理删除}，只打删除标记，误删可以恢复。
  \item \textbf{审计日志}：所有\textbf{写操作}都会记录到审计日志（谁、在什么时间、对什么资源、做了什么）。
  \item \textbf{自我保护}：多处关键任务带「熔断」——当发现异常范围过大（可能是配置错了）时，宁可什么都不做也不批量删数据。
\end{itemize}

% =====================================================================
\chapter{部署与运行}
% =====================================================================

\begin{introduction}
  \item 需要准备什么环境
  \item 怎么构建出可部署的程序
  \item 怎么在 Linux 服务器上跑起来
\end{introduction}

\section{环境要求}

\begin{itemize}
  \item 运行环境：\textbf{Java 21}、\textbf{MySQL 8}。
  \item 构建环境（仅打包时需要）：Node.js 20+、Zig 0.16.x、CMake 3.24+、Ninja、NASM。
  \item Linux 运行还需要：glibc 2.28+、\code{libstdc++.so.6}（含 \code{GLIBCXX\_3.4.22}）、\code{libgcc\_s.so.1}。
\end{itemize}

\begin{notebox}[部署机不需要构建工具]
  构建（打包 JAR）和运行是两回事。JAR 一旦打好，\textbf{部署服务器上只要有 Java 21 即可}，
  不需要 Node、Maven、Zig，也不需要单独部署前端。
\end{notebox}

\section{构建可部署的 JAR}

在开发机（已装好构建环境）上执行：

\begin{stepbox}[打包成一个 JAR]
\begin{verbatim}
cd backend
.\mvnw.cmd clean package     # Windows
./mvnw clean package         # Linux / macOS
\end{verbatim}
\end{stepbox}

这条命令会自动安装隔离的 Node、构建 React 前端、编译两套原生图片组件（Windows / Linux），
并生成 \code{backend/target/photolib-backend-0.1.0-SNAPSHOT.jar}。这个 JAR 里已经包含了后端、前端和图片组件。

\begin{tipbox}[用 GitHub Actions 自动发版]
  项目已配置好自动打包流水线。发布一个版本只需要推一个标签：

  \begin{verbatim}git tag v1.0.0 && git push origin v1.0.0\end{verbatim}

  流水线会跑测试、打包、并自动创建带附件的 Release。带连字符的标签（如 \code{v1.2.0-rc.1}）会当作预发布处理。
\end{tipbox}

\section{在 Linux 上用 systemd 运行}

建议用专用用户和 \code{/opt/photolib} 目录部署，把 JAR 和生产 \code{.env} 放进去。一个最简 systemd 单元如下：

\begin{verbatim}
[Unit]
Description=PhotoLib
After=network-online.target
Wants=network-online.target

[Service]
Type=simple
User=photolib
Group=photolib
WorkingDirectory=/opt/photolib
ExecStart=/usr/bin/java -jar /opt/photolib/photolib.jar
Restart=on-failure
RestartSec=5
SuccessExitStatus=143

[Install]
WantedBy=multi-user.target
\end{verbatim}

启动并检查健康状态：

\begin{verbatim}
sudo systemctl daemon-reload
sudo systemctl enable --now photolib
sudo systemctl status photolib
curl --fail http://127.0.0.1:8080/api/v1/actuator/health
\end{verbatim}

\begin{warnbox}[务必放在 HTTPS 反向代理后面]
  应用本身跑在 8080 端口。生产环境必须在它前面加一层 HTTPS 反向代理（如 Nginx），
  并且\textbf{只允许反向代理访问应用端口}，不要把 8080 直接暴露到公网。
\end{warnbox}

% =====================================================================
\chapter{生产环境配置}
% =====================================================================

\begin{introduction}
  \item 敏感配置怎么注入
  \item 有哪些必须关注的环境变量
\end{introduction}

\section{配置从哪来}

敏感配置（数据库密码、存储密钥等）只能通过\textbf{环境变量}、JAR 工作目录下受限权限的 \code{.env}，
或密钥管理服务注入。\textbf{绝对不要}把 AccessKey、桶地址、签名密钥或 \code{.env} 提交到代码仓库。

\section{常用环境变量}

\begin{longtable}{L{5.4cm} L{7.1cm}}
  \toprule
  \textbf{环境变量} & \textbf{说明} \\
  \midrule
  \endfirsthead
  \toprule
  \textbf{环境变量} & \textbf{说明} \\
  \midrule
  \endhead
  \code{DB\_URL / DB\_USERNAME / DB\_PASSWORD} & MySQL 连接配置 \\
  \code{ADMIN\_INITIAL\_PASSWORD} & 首次启动的管理员密码，必须用强随机值 \\
  \code{AUTH\_SECURE\_COOKIE} & 生产环境必须为 \code{true} \\
  \code{STORAGE\_MODE} & 本地开发填 \code{local}，生产填 \code{oss} \\
  \code{OSS\_BUCKET / OSS\_ENDPOINT / OSS\_PUBLIC\_ENDPOINT} & 私有桶、服务端 Endpoint、浏览器访问 Endpoint \\
  \code{OSS\_ACCESS\_KEY\_ID / \_SECRET} & 最小权限的 RAM 用户凭据 \\
  \code{OSS\_CORS\_ALLOWED\_ORIGINS} & 允许浏览器直传的站点来源，生产应显式配置 \\
  \code{OSS\_SIGNATURE\_WINDOW} & 预览地址签发时刻取整窗口，默认 \code{5m}，须小于下载 TTL \\
  \code{PREVIEW\_COMPRESSION\_RATIO} & 预览图质量，0 到 1，默认 \code{0.6} \\
  \code{LOGIN\_MAX\_IDENTIFIER\_FAILURES} & 单账号连续失败几次后锁定，默认 \code{5} \\
  \code{LOGIN\_MAX\_ADDRESS\_FAILURES} & 单客户端地址累计失败几次后锁定，默认 \code{40} \\
  \code{LOGIN\_FAILURE\_WINDOW / LOGIN\_LOCK\_DURATION} & 失败计数窗口 / 锁定时长，均默认 \code{15m} \\
  \code{PHOTO\_PROCESSING\_THREADS} & 图片处理线程数，1 到 32，默认 \code{1} \\
  \code{STARTUP\_MISSING\_OBJECT\_CLEANUP\_ENABLED} & 启动时是否清理「成品图已丢失」的记录，默认 \code{true} \\
  \code{DATABASE\_BACKUP\_ENABLED} & 是否每天 0 点自动备份数据库，默认 \code{true} \\
  \code{DATABASE\_BACKUP\_RETENTION} & 备份保留时长，默认 \code{30d} \\
  \code{DATABASE\_BACKUP\_MINIMUM\_RETAINED} & 无论保留期，至少保留几份成功备份，默认 \code{7} \\
  \code{DIRECTMAIL\_*} & 邮件服务的地域、发信地址与凭据 \\
  \bottomrule
\end{longtable}

\begin{warnbox}[首次上线务必做的三件事]
  \begin{enumerate}
    \item 用强随机值设置 \code{ADMIN\_INITIAL\_PASSWORD}，首次登录后立刻改密。
    \item 生产环境不要使用 \code{local} Profile、默认签名密钥或默认密码。
    \item 把 \code{OSS\_CORS\_ALLOWED\_ORIGINS} 收窄到你真实的站点来源，不要长期用 \code{*}。
  \end{enumerate}
\end{warnbox}

% =====================================================================
\chapter{账号与权限体系}
% =====================================================================

\begin{introduction}
  \item 三种角色分别管什么
  \item 权限组、权限码、数据范围是怎么配合的
  \item 「图库可见范围」这个容易被忽略的开关
\end{introduction}

\section{三种角色}

\begin{table}[htbp]
  \centering
  \caption{系统的三种角色}
  \begin{tabular}{L{2.4cm} L{2.2cm} L{7.4cm}}
    \toprule
    \textbf{角色} & \textbf{代码} & \textbf{默认职责} \\
    \midrule
    管理员 & \code{ADMIN} & 管理账号、权限组、校区及全部业务资源，查看审计与告警 \\
    部长 & \code{MINISTER} & 管理项目和需求，维护采用记录，确认工时并导出统计 \\
    校区负责人 & \code{CAMPUS\_MANAGER} & 处理授权校区内的需求、图片、通讯录与工时 \\
    \bottomrule
  \end{tabular}
\end{table}

系统\textbf{不开放注册}。首次启动时只有一个管理员账号，其余账号都由管理员创建。

\section{权限组与权限码}

角色只是一个大类，真正决定「能点哪些菜单、做哪些操作」的是\textbf{权限组}。每个账号属于一个权限组，
权限组里勾选了一批\textbf{权限码}。内置权限组（\code{ADMIN} / \code{MINISTER} / \code{CAMPUS\_MANAGER}）
的名称和数据范围不可改，但你可以\textbf{新建自定义权限组}，自由组合权限码。

\begin{table}[htbp]
  \centering
  \caption{主要权限码（在「系统管理 → 权限组」里按分类勾选）}
  \begin{tabular}{L{4.8cm} L{7.6cm}}
    \toprule
    \textbf{权限码} & \textbf{含义} \\
    \midrule
    \code{PROJECT\_VIEW / \_CREATE / \_COMPLETE} & 选题的查看 / 新建编辑发布 / 标记完成 \\
    \code{PROJECT\_ADOPT} & 标记或取消图片被引（采用） \\
    \code{REQUEST\_VIEW / \_CREATE / \_CONFIRM} & 需求的接受提交 / 新建发布 / 确认退回 \\
    \code{REQUEST\_PHOTO\_MANAGE} & 管理需求图片（上传、下载、删除） \\
    \code{PHOTO\_VIEW / \_UPLOAD / \_DOWNLOAD / \_DELETE} & 图库的看 / 传 / 下 / 删 \\
    \code{WORKLOG\_SUBMIT / \_CONFIRM / \_EXPORT} & 工时的填报 / 确认退回 / 导出 \\
    \code{DIRECTORY\_VIEW / \_MANAGE} & 通讯录的查看 / 维护 \\
    \code{STATISTICS\_DOWNLOAD} & 统计数据的查看与下载 \\
    \code{RECRUITMENT\_VIEW / \_PUBLISH} & 招募的查看 / 创建发布 \\
    \code{FEATURED\_MANAGE / DOC\_MANAGE / MESSAGE\_SEND} & 好图精选 / 文档 / 消息 的管理发送 \\
    \code{MANAGER\_CAMPUS\_ASSIGN} & 重新指定校区负责人的授权校区 \\
    \bottomrule
  \end{tabular}
\end{table}

\section{数据范围与图库可见范围}

除了「能做什么」，权限组还控制「能看到谁的数据」。这里有两个正交（互不影响）的维度：

\begin{itemize}
  \item \textbf{数据范围（data\_scope）}：管的是\textbf{校区授权}。\code{GLOBAL} 看全站，校区范围只看授权校区。
  \item \textbf{图库可见范围（photo\_visibility）}：管的是\textbf{图库里看得到谁的图}。\code{SELF} 只看本人上传、
    \code{CAMPUS} 看授权校区内全部、\code{GLOBAL} 看全站。
\end{itemize}

\begin{tipbox}[一个常见的调整场景]
  校区负责人默认图库可见范围是 \code{SELF}（只能看自己上传的图）。如果你希望某校区的负责人\textbf{互相看得到}
  本校区其他人上传的图片，就把他所在权限组的图库可见范围改成 \code{CAMPUS}。
  注意：放宽「看得到」不会连带放开「改、删、标记采用」——那些始终限本人上传。
\end{tipbox}

% =====================================================================
\chapter{后台管理}
% =====================================================================

\begin{introduction}
  \item 「系统管理」页面能做哪些事
  \item 每一块的用途和注意点
\end{introduction}

管理员登录后，左侧菜单底部有\menu{系统管理}入口（其他角色看不到）。它把管理能力集中在几个面板里。

\section{账号管理}

创建、编辑、停用、删除账号，重置密码，指定所属权限组。要点：

\begin{itemize}
  \item 新账号首次登录会\textbf{强制改密}。
  \item 邮箱可留空；若填写，在未删除用户里必须唯一（保存时会自动转小写）。
  \item 登录标识可以是用户名，\textbf{也可以是邮箱}。
  \item 删除账号采用「释放用户名后软删除」：历史记录仍在，同一个用户名日后可以重新创建。
  \item \textbf{不能删除当前登录账号，也不能删除唯一仍启用的管理员。}
  \item 停用、改密、删除账号都会\textbf{撤销该用户的登录会话}。
\end{itemize}

\section{校区与负责人校区}

维护校区列表（启用/停用），以及给校区负责人\textbf{指定授权校区}。校区停用后，其下的需求就不能再被接受。

\section{权限组}

新建自定义权限组、勾选权限码、设置数据范围和图库可见范围。内置三组不可删、名称与数据范围不可改，
但图库可见范围可以调整（见上一章）。

\section{品牌定制}

配置产品名、Slogan 和图标（内置图标或上传图片）。保存后侧栏会立即更新，登录页、招募页等也会跟随。
前端所有产品标识都从这里读取，\textbf{不再硬编码任何产品名}。

\section{审计日志}

\begin{itemize}
  \item 仅管理员可访问。可按操作人、HTTP 动作、资源类型、关键词、起止日期筛选并分页。
  \item 结束日期按「次日零点之前」处理，所以选中的结束日\textbf{包含一整天}。
  \item 可\textbf{导出 CSV}（带 BOM，Excel 打开中文不乱码）。当前单次导出上限 10 万条。
  \item 审计只记录\textbf{写操作}；登录尝试（含失败）也会记录，但不会记录密码或请求体。
\end{itemize}

\section{数据备份}

\menu{系统管理 → 数据备份}，这是\textbf{只对管理员开放}的能力，没有对应权限码、也无法授予自定义权限组。详见下一章。

\begin{notebox}[存储对账与告警]
  系统会定期核对「数据库里的图片记录」和「对象存储里的实际文件」是否对得上，发现异常时写入\textbf{管理员告警}。
  你会在后台看到这些告警。它们大多是提醒你人工确认，而不是系统已经删了数据——恰恰相反，系统在拿不准时会保守地\textbf{不删}。
\end{notebox}

% =====================================================================
\chapter{数据安全与运维}
% =====================================================================

\begin{introduction}
  \item 数据库备份和回滚怎么用
  \item 系统有哪些自动的自我保护任务
  \item 匿名报名入口的运维红线
\end{introduction}

\section{数据库备份与回滚}

系统每天凌晨 0 点（Asia/Shanghai）自动把整库业务数据备份到对象存储。管理员也可以手动备份、下载、
导入外部备份文件，或把数据库\textbf{回滚}到某个备份。要点：

\begin{itemize}
  \item 备份是\textbf{逻辑备份}（gzip 压缩的 JSON，只含数据不含表结构），不依赖服务器上装 mysqldump。
  \item 表结构由 Flyway 管理，所以\textbf{只能回滚到与当前结构版本一致的备份}，跨版本的备份界面上不可回滚。
  \item 回滚前会自动生成一份 \code{PRE\_RESTORE} 兜底备份，误操作后还能再退回来。
  \item 回滚在\textbf{单个事务}里完成，中途失败整体回滚，不会留下「删了一半」的库。
  \item 回滚\textbf{只替换数据库}，不会删除对象存储里的图片文件；但回滚后登录会话可能全部失效。
\end{itemize}

\begin{warnbox}[单实例假设]
  备份/回滚的并发控制目前\textbf{假设单实例部署}。如果将来要多实例部署，必须先改成数据库层的抢占锁，
  否则可能出现两个实例同时回滚。
\end{warnbox}

\section{自动的自我保护任务}

系统在每次启动和运行期间会跑几个自动任务，理解它们有助于你看懂日志和告警：

\begin{table}[htbp]
  \centering
  \caption{几个关键的自动任务}
  \begin{tabular}{L{4.6cm} L{7.8cm}}
    \toprule
    \textbf{任务} & \textbf{作用与保护} \\
    \midrule
    启动缺失对象清理 & 核对成品图是否还在存储里；\textbf{只有精确确认对象不存在}才软删记录。已被采用的图只告警不删；缺失比例过高（疑似配置错误）时一张都不删并告警。 \\
    原图清理 & 到保留期且已生成可用图片的原图才会被删。 \\
    预览图重建 & 修改压缩比或换编码器后，在后台逐张重建预览图，不阻塞启动，失败只影响单张。 \\
    存储对账 & 定期核对数据库与存储，异常写管理员告警，不据此批量删数据。 \\
    \bottomrule
  \end{tabular}
\end{table}

\begin{tipbox}[看到告警不要慌]
  这些任务的设计原则是「宁可不动，也不误删」。看到带 ABORTED 字样的熔断告警，
  通常意味着\textbf{存储配置出了问题}（桶名、Endpoint、前缀或挂载点配错），应先排查配置，而不是怀疑数据真丢了。
  （告警名形如 \texttt{PHOTO\_MISSING\_OBJECT\_CLEANUP\_ABORTED}。）
\end{tipbox}

\section{登录防护}

应用内做了两层防护：按\textbf{账号}和按\textbf{客户端地址}分别计失败次数并锁定；登录尝试全量进审计。
但这\textbf{不替代网关限流}——应用只能看到已经到达它的请求。请务必在反向代理/网关上为 \code{/api/v1/auth/login}
配置独立限流。

\section{成员招募的匿名入口}

成员招募的公开报名页（\code{/\#/recruitment}）是系统里\textbf{唯一对未登录访客开放写操作}的功能。
上线前请注意：

\begin{warnbox}[网关必须承担的防护]
  \code{/api/v1/public/recruitments/**} 完全匿名。应用内的限流只是\textbf{有界的纵深防护，不能替代网关限流}：
  部署在反向代理后面时，应用拿到的往往是代理 IP，会被当作共享地址\textbf{直接放行}。
  因此\textbf{必须}在网关上为这些路径配置分布式限流。开放招募前，请用真实大小的样本包实测一次上传是否顺畅。
\end{warnbox}

% =====================================================================
\chapter{旧系统数据迁移（如需）}
% =====================================================================

如果要从旧的 PhotoWarehouse 系统（PostgreSQL + OSS）迁移数据，脚本在 \code{scripts/} 目录下，采用可重入的两阶段流程：

\begin{enumerate}
  \item \code{migrate\_photowarehouse.py export}：从旧库导出完整 JSON 包。
  \item \code{migrate\_photowarehouse.py import}：复制 OSS 对象、写入新库、归档旧行。
  \item \code{restore\_project\_photos.py}：按旧标签规则重建「一图多项目」的相册归属。
\end{enumerate}

\begin{warnbox}[迁移前后须知]
  \begin{itemize}
    \item 执行前\textbf{必须备份}数据库与 Bucket，并先启动一次新后端让 Flyway 升级到最新版本。
    \item 图片对象\textbf{不下载、不重传}，只复用旧的存储 key。
    \item 旧数据无法可靠推断的字段一律留空，尤其\textbf{不能猜测图片校区}。
    \item 确认迁移完成后，务必关闭迁移开关 \code{IS\_MIGRATE}。
  \end{itemize}
\end{warnbox}

详细字段映射与命令见项目 \code{docs/} 下的迁移文档。

% =====================================================================
\chapter{日常巡检清单}
% =====================================================================

\begin{stepbox}[建议定期检查]
  \begin{enumerate}
    \item 健康检查地址 \code{/api/v1/actuator/health} 是否正常。
    \item 后台\textbf{管理员告警}里有没有新的存储/备份/迁移告警。
    \item 数据库\textbf{每日备份}是否成功（可在数据备份面板查看）。
    \item 审计日志量是否过大，是否需要制定保留与清理策略（登录审计会显著增加写入量）。
    \item 反向代理/网关的\textbf{限流规则}是否仍对登录和招募匿名接口生效。
    \item 存储桶的 RAM 凭据、CORS 配置是否仍然正确、最小权限。
  \end{enumerate}
\end{stepbox}

\begin{notebox}[更多技术细节]
  本指南是快速上手向导。更深入的模块业务规则、不变量和「改动会破坏什么」，
  请查阅项目根目录的 \code{AGENTS.md} 与 \code{README.md}。
\end{notebox}

\end{document}
